% ============================================================
% CONCLUSION - target: 150--300 words
% ============================================================

\section{Conclusion}

This paper presents a systematic, cross-family evaluation of chain-of-thought
faithfulness in open-weight reasoning models, testing 12 models from 9
architectural families across 41,832 inference runs with six categories of
reasoning hints. Faithfulness rates range from 39.7\% to 89.9\% across 12
models from 9 families, with social-pressure hints (consistency 35.5\%,
sycophancy 53.9\%) proving hardest to acknowledge. Keyword-based analysis
indicates that thinking tokens are acknowledged far more frequently
(approximately 87.5\%) than answer text (approximately 28.6\%).

These findings carry both theoretical and practical implications. On the
theoretical side, training methodology shapes faithfulness more than model
scale, and keyword-based analysis suggests that models develop an internal
capacity to detect hint influence (thinking-token acknowledgment averaging
approximately 87.5\%) that is systematically suppressed in their externally
visible outputs. On the practical
side, safety teams should (1)~prioritize thinking-token monitoring over
answer-text monitoring, since answer text captures less than a third of
available faithfulness signal; (2)~select models with demonstrated faithfulness
for safety-critical deployments, noting the wide variation across families; and
(3)~treat published faithfulness numbers as classifier-dependent estimates rather
than absolute measurements, given the 12.9-percentage-point gap observed
between the two classifier pipelines~\cite{young2026classifier}.

These findings are subject to the caveats of API-based evaluation,
multiple-choice format, conservative classifiers, and the documented sensitivity
of faithfulness rates to classifier methodology.

Looking forward, three directions are particularly promising: (1)~extending
faithfulness evaluation to open-ended generation tasks where models have more
freedom to construct alternative justifications; (2)~combining surface-level
CoT analysis with mechanistic interpretability techniques to detect faithfulness
gaps that are invisible to text-based classifiers; and (3)~developing
training-time interventions that improve faithfulness without sacrificing task
performance. As reasoning models are increasingly deployed in high-stakes
settings, the gap between what models think and what they say represents a
fundamental challenge for CoT-based safety monitoring, one that demands both
better measurement tools and training approaches that align internal reasoning
with external communication.
